\chapter{ТЕСТИРОВАНИЕ РАЗРАБОТАННОГО ВЕБ-ПРИЛОЖЕНИЯ}

После завершения разработки веб-приложения «shasl mp» необходимо провести комплексное тестирование, подтверждающее его работоспособность, надёжность и безопасность. Цель тестирования — выявить возможные ошибки в реализации ключевых модулей, проверить корректность работы всех функциональных требований и убедиться, что система устойчива к некорректным входным данным и злоумышленным действиям.

В ходе тестирования решаются следующие задачи:

\begin{enumerate}
\item Проверка функциональности подсистемы авторизации (регистрация, вход, защита маршрутов, выход).
\item Проверка интеграции с API маркетплейсов (создание отчёта, проверка статуса, скачивание).
\item Проверка корректности расчётов модулей аналитики (юнит-экономика, кластерный анализ, визуализация продаж).
\item Проверка обработки некорректных форматов файлов и ошибочных сценариев.
\item Проверка защищённости системы от SQL-инъекций, XSS-атак и несанкционированного доступа.
\end{enumerate}

Тестирование проводилось в среде разработки (Windows 11, Node.js 20, браузер Google Chrome) с использованием реальных отчётов, выгруженных из личного кабинета Ozon, а также сгенерированных тестовых данных. Для каждого тест-кейса описан ожидаемый результат, фактические результаты и зафиксированы отклонения (при наличии). По итогам тестирования сделан вывод о готовности системы к эксплуатации.

Виды тестирования, применённые в работе (см. табл.~\ref{tab:testing-types}).

\begin{table}[H]
\caption{Виды тестирования}
\label{tab:testing-types}
\begin{tabular}{|p{4.5cm}|p{10.5cm}|}
\hline
Вид тестирования & Что проверяет \\ \hline
Функциональное & Соответствие реализованных функций требованиям (авторизация, API, расчёты) \\ \hline
Стресс-тестирование & Поведение системы при некорректных входных данных (неверный формат файла, пустой файл) \\ \hline
Тестирование безопасности & Защищённость от SQL-инъекций, XSS-атак, подделки сессий \\ \hline
\end{tabular}
\end{table}

\section{Тестирование подсистемы авторизации}

Подсистема авторизации является критически важным компонентом, так как она разграничивает доступ к аналитическим данным разных пользователей. Тестирование проводилось по следующим сценариям: регистрация нового пользователя (см. табл.~\ref{tab:testReg}), вход с валидными и невалидными данными (см. табл.~\ref{tab:testEnter}), доступ к защищённым страницам (см. табл.~\ref{tab:testAuth}), поведение при активной сессии (см. табл.~\ref{tab:testSession}) и выход из системы (см. табл.~\ref{tab:testExit}).

Среда тестирования авторизации:
\begin{itemize}
\item браузер: Google Chrome (версия 124);
\item сервер: Node.js 20, локальный хост (http://localhost:3000);
\item база данных: SQLite (файл users.db);
\item инструменты: консоль разработчика (F12) для проверки cookies и сетевых запросов.
\end{itemize}

\begin{table}[H]
\caption{\label{tab:testReg}Тестирование регистрации пользователя}
\small
\begin{tabular}{|p{4cm}|p{4cm}|p{3cm}|p{4cm}|}
\hline
Действие & Ожидаемый результат & Фактический результат & Результат \\ \hline
Ввод email: test@example.com, пароль: 12345678 & Пользователь создан, перенаправление на /login, сообщение об успехе & Соответствует & Положительный \\ \hline
Ввод email: test@example.com (существующий), пароль: 12345678 & Ошибка: «Пользователь с таким email уже существует» & Соответствует & Положительный \\ \hline
Ввод email: test@example.com, пароль: 123 (<6 символов) & Ошибка: «Пароль должен быть не менее 6 символов» & Соответствует & Положительный \\ \hline
Ввод некорректного email: testexample.com & Ошибка валидации на клиенте (HTML5 required) & Соответствует & Положительный \\ \hline
\end{tabular}
\end{table}

\begin{table}[H]
\caption{\label{tab:testEnter}Тестирование входа пользователя}
\small
\begin{tabular}{|p{4cm}|p{4cm}|p{3cm}|p{4cm}|}
\hline
Действие & Ожидаемый результат & Фактический результат & Результат \\ \hline
Вход с test@example.com / 12345678 (существующий) & JWT-токен создан, установлена HTTP-only cookie, перенаправление на / & Соответствует & Положительный \\ \hline
Вход с неверным паролем & Ошибка: «Неверный email или пароль», cookie не устанавливается & Соответствует & Положительный \\ \hline
Вход с несуществующим email & Ошибка: «Неверный email или пароль» & Соответствует & Положительный \\ \hline
\end{tabular}
\end{table}

\begin{table}[H]
\caption{\label{tab:testAuth}Тестирование доступа без авторизации}
\small
\begin{tabular}{|p{4cm}|p{4cm}|p{3cm}|p{4cm}|}
\hline
Действие & Ожидаемый результат & Фактический результат & Результат \\ \hline
Попытка открыть / (главная страница) без авторизации & Перенаправление на /login & Соответствует & Положительный \\ \hline
Попытка открыть страницу /ozon без авторизации & Перенаправление на /login & Соответствует & Положительный \\ \hline
Попытка открыть API-эндпоинт без токена & Перенаправление на /login & Соответствует & Положительный \\ \hline
\end{tabular}
\end{table}

\begin{table}[H]
\caption{\label{tab:testSession}Тестирование поведения при активной сессии}
\small
\begin{tabular}{|p{4cm}|p{4cm}|p{3cm}|p{4cm}|}
\hline
Действие & Ожидаемый результат & Фактический результат & Результат \\ \hline
Авторизованный пользователь пытается открыть /login & Перенаправление на / & Соответствует & Положительный \\ \hline
Авторизованный пользователь пытается открыть /register & Перенаправление на / & Соответствует & Положительный \\ \hline
\end{tabular}
\end{table}

\begin{table}[H]
\caption{\label{tab:testExit}Тестирование выхода из системы}
\small
\begin{tabular}{|p{4cm}|p{4cm}|p{3cm}|p{4cm}|}
\hline
Действие & Ожидаемый результат & Фактический результат & Результат \\ \hline
Нажатие кнопки «Выйти» & Cookie session удалена, перенаправление на /login & Соответствует & Положительный \\ \hline
\end{tabular}
\end{table}

Результаты тестирования (см. табл.~\ref{tab:testReg}--\ref{tab:testExit}) авторизации: все тест-кейсы пройдены успешно. Подсистема регистрирует новых пользователей, корректно проверяет пароли, защищает маршруты от неавторизованного доступа, правильно обрабатывает ошибки и обеспечивает безопасное хранение токенов в HTTP-only cookies.

\section{Тестирование скачивания отчёта через API}

Для автоматизации аналитики необходимо, чтобы приложение корректно взаимодействовало с API маркетплейса Ozon. Тестирование охватывает три этапа: создание отчёта POST /v1/report/sales (см. табл.~\ref{tab:testCreateReport}), проверку готовности POST /v1/report/info (см. табл.~\ref{tab:testWaitForReport}) и скачивание файла по полученной ссылке (см. табл.~\ref{tab:testDownloadReport}). Также проверяются ошибочные сценарии: неверные параметры запроса, недействительные API-ключи и превышение времени ожидания.

Среда тестирования:
\begin{itemize}
\item API-эндпоинт: https://api-seller.ozon.ru/v1/;
\item тестовый Client ID и API-ключ (реальные, из личного кабинета);
\item инструменты: консоль разработчика (Network), логи сервера, Postman для верификации.
\end{itemize}

\begin{table}[!htb]
\caption{\label{tab:testCreateReport}Тестирование создания отчёта (createReport)}
\small
\begin{tabular}{|p{4cm}|p{4cm}|p{3cm}|p{4cm}|}
\hline
Действие & Ожидаемый результат & Фактический результат & Результат \\ \hline
Запрос отчёта за корректный период (например, последние 7 дней) & Сервер возвращает HTTP 200, в теле ответа поле code (уникальный идентификатор отчёта) & Соответствует & Положительный \\ \hline
Запрос с датой начала больше даты окончания & Ошибка API (400), сообщение об ошибке валидации & Соответствует & Положительный \\ \hline
Запрос с неверным Client ID & Ошибка авторизации (401) & Соответствует & Положительный \\ \hline
Запрос с неверным API-ключом & Ошибка авторизации (401) & Соответствует & Положительный \\ \hline
\end{tabular}
\end{table}

\begin{table}[!htb]
\caption{\label{tab:testWaitForReport}Тестирование проверки готовности отчёта (waitForReport)}
\small
\begin{tabular}{|p{4cm}|p{4cm}|p{3cm}|p{4cm}|}
\hline
Действие & Ожидаемый результат & Фактический результат & Результат \\ \hline
Опрос статуса через 5 секунд после создания (статус pending) & Возвращается status: "pending", цикл продолжается & Соответствует & Положительный \\ \hline
Опрос после завершения формирования (статус success) & Возвращается status: "success" и поле file (URL для скачивания) & Соответствует & Положительный \\ \hline
Опрос с несуществующим code & Ошибка API (404), цикл прерывается & Соответствует & Положительный \\ \hline
Превышение максимального времени ожидания (10 попыток по 5 секунд) & Ошибка «Превышено время ожидания отчёта» & Соответствует & Положительный \\ \hline
\end{tabular}
\end{table}


\begin{table}[!htb]
\caption{\label{tab:testDownloadReport}Тестирование скачивания отчёта (downloadReport)}
\small
\begin{tabular}{|p{4cm}|p{4cm}|p{3cm}|p{4cm}|}
\hline
Действие & Ожидаемый результат & Фактический результат & Результат \\ \hline
GET-запрос по URL, полученному из waitForReport & Файл (CSV/XLSX) успешно загружен, сохранён на сервере & Соответствует & Положительный \\ \hline
Скачивание файла в кодировке Windows-1251 & Автоматическая конвертация в UTF-8, содержимое читаемо & Соответствует & Положительный \\ \hline
Повторное использование одноразовой ссылки & Ошибка 404 (файл не найден) & Соответствует & Положительный \\ \hline
\end{tabular}
\end{table}

Результаты тестирования API-взаимодействия (см. табл.~\ref{tab:testCreateReport}--\ref{tab:testDownloadReport}): все тест-кейсы пройдены успешно. Приложение корректно создаёт отчёты, обрабатывает цикл ожидания, скачивает файлы, преобразует кодировки и выдаёт понятные ошибки при неверных параметрах или проблемах авторизации. Время получения отчёта за 30 дней составляет в среднем 15–20 секунд (с учётом ожидания формирования на стороне Ozon).

\section{Тестирование модуля юнит-экономики}

Модуль юнит-экономики отвечает за расчёт ключевых финансовых метрик: выручки, себестоимости, прибыли с единицы товара и общей прибыли за период. Корректность этих расчётов критически важна для принятия управленческих решений продавцом. Тестирование проводилось на основе реальных отчётов о продажах Ozon, а также на синтетических данных с известными значениями.

Среда тестирования:
\begin{itemize}
\item исходные данные: CSV-отчёт о продажах (заказы за период 01.04.2025 – 30.04.2025);
\item эталонные значения: рассчитаны вручную в Excel;
\item инструменты: логи приложения, вывод модуля в консоль.
\end{itemize}

Входные параметры модуля (см. табл.~\ref{tab:testUnitEcon}):
\begin{table}[htb]
\caption{\label{tab:testUnitEcon}Входные параметры модуля юнит-экономики}
\small
\begin{center}
\vspace{-0.5cm}\begin{tabular}{|l|c|}
\hline
Параметр & Значение (для тестового набора) \\ \hline
Цена товара & 1000 руб. \\ \hline
Количество в заказе & 1–5 шт. \\ \hline
Себестоимость закупки & 400 руб. \\ \hline
Комиссия маркетплейса & 15\% от цены \\ \hline
Логистика & 50 руб. за единицу \\ \hline
Статус заказа & Доставлен / Отменён \\ \hline
\end{tabular}
\end{center}
\end{table}

Формулы расчёта, реализованные в модуле (см. табл.~\ref{tab:equation}):
\begin{table}[htb]
\caption{\label{tab:equation}Формулы расчёта метрик}
\small
\begin{center}
\vspace{-0.5cm}\begin{tabular}{|l|l|}
\hline
Метрика & Формула \\ \hline
Выручка & Сумма (цена $\times$ количество) по доставленным заказам \\ \hline
Себестоимость единицы & Закупка + (цена $\times$ комиссия) + логистика \\ \hline
Прибыль с единицы & Цена – Себестоимость единицы \\ \hline
Общая прибыль & Сумма (прибыль с единицы $\times$ количество) \\ \hline
\end{tabular}
\end{center}
\end{table}

Расчёт выручки товара по заданным исходным данным представлен в таблице~\ref{tab:testEquation}). Расчёт себестоимости единицы товара по заданным исходным данным протестирован и представлен в таблице~\ref{tab:testCostOfProduction}). Расчёт прибыли с единицы товара по заданным исходным данным представлен в таблице~\ref{tab:testPerUnit}. Тестирование расчёта общей прибыли за период по заданным исходным данным представлен в таблице~\ref{tab:testCalProf}.
\begin{table}[!htb]
\caption{\label{tab:testEquation}Тестирование расчёта выручки}
\small
\begin{center}
\vspace{-0.5cm}\begin{tabular}{|p{4cm}|p{3cm}|p{3cm}|p{4cm}|}
\hline
Исходные данные & Ожидаемая выручка & Фактическая выручка & Результат \\ \hline
2 заказа по 1000 руб. (доставлены), 1 отменённый на 1500 руб. & 2000 руб. & 2000 руб. & Положительный \\ \hline
5 заказов: 3 по 1000 руб., 2 по 500 руб. (все доставлены) & 3000 + 1000 = 4000 руб. & 4000 руб. & Положительный \\ \hline
Пустой отчёт (нет доставленных заказов) & 0 руб. & 0 руб. & Положительный \\ \hline
\end{tabular}
\end{center}
\end{table}

\begin{table}[!htb]
\caption{\label{tab:testCostOfProduction}Тестирование расчёта себестоимости единицы}
\small
\begin{center}
\vspace{-0.5cm}\begin{tabular}{|p{5cm}|p{3cm}|p{3cm}|p{4cm}|}
\hline
Исходные данные & Ожидаемая себестоимость & Фактическая себестоимость & Результат \\ \hline
Цена 1000 руб., закупка 400 руб., комиссия 15\% (150 руб.), логистика 50 руб. & 400 + 150 + 50 = 600 руб. & 600 руб. & Положительный \\ \hline
Цена 2000 руб., закупка 800 руб., комиссия 15\% (300 руб.), логистика 50 руб. & 800 + 300 + 50 = 1150 руб. & 1150 руб. & Положительный \\ \hline
Цена 500 руб., закупка 200 руб., комиссия 15\% (75 руб.), логистика 50 руб. & 200 + 75 + 50 = 325 руб. & 325 руб. & Положительный \\ \hline
\end{tabular}
\end{center}
\end{table}

\begin{table}[htb]
\caption{\label{tab:testPerUnit}Тестирование расчёта прибыли с единицы}
\small
\begin{center}
\vspace{-0.5cm}\begin{tabular}{|p{7cm}|p{2.5cm}|p{2.5cm}|p{3cm}|}
\hline
Исходные данные & Ожидаемая прибыль & Фактическая прибыль & Результат \\ \hline
Цена 1000 руб., себестоимость 600 руб. & 400 руб. & 400 руб. & Положительный \\ \hline
Цена 2000 руб., себестоимость 1150 руб. & 850 руб. & 850 руб. & Положительный \\ \hline
Цена 500 руб., себестоимость 325 руб. & 175 руб. & 175 руб. & Положительный \\ \hline
\end{tabular}
\end{center}
\end{table}

\begin{table}[htb]
\caption{\label{tab:testCalProf}Тестирование расчёта общей прибыли за период}
\small
\begin{center}
\vspace{-0.5cm}\begin{tabular}{|p{7cm}|p{2.5cm}|p{2.5cm}|p{3cm}|}
\hline
Исходные данные & Ожидаемая прибыль & Фактическая прибыль & Результат \\ \hline
2 единицы с прибылью 400 руб., 1 единица с прибылью 175 руб. & $400 \times 2 + 175 = 975$ руб. & 975 руб. & Положительный \\ \hline
10 единиц с прибылью 400 руб., 5 единиц с прибылью 175 руб. & $4000 + 875 = 4875$ руб. & 4875 руб. & Положительный \\ \hline
Отменённые заказы (не включаются в расчёт) & 0 руб. & 0 руб. & Положительный \\ \hline
\end{tabular}
\end{center}
\end{table}

Обработка ошибок и пограничных случаев по заданным сценариям (см. табл.~\ref{tab:testError})

\begin{table}[H]
\caption{\label{tab:testError}Тестирование обработки ошибок}
\small
\begin{center}
\vspace{-0.5cm}\begin{tabular}{|p{3cm}|p{6cm}|p{3cm}|p{3cm}|}
\hline
Сценарий & Ожидаемый результат & Фактический результат & Результат \\ \hline
Отсутствует файл отчёта & Ошибка: «Файл не найден» & Соответствует & Положительный \\ \hline
Пустая цена (null) в отчёте & Пропуск строки, логирование предупреждения & Соответствует & Положительный \\ \hline
Отрицательное количество & Пропуск строки, логирование предупреждения & Соответствует & Положительный \\ \hline
\end{tabular}
\end{center}
\end{table}

Результаты тестирования модуля юнит-экономики (см. табл.~\ref{tab:testEquation}--\ref{tab:testError}): все тест-кейсы пройдены успешно. Расчёт выручки, себестоимости и прибыли соответствует эталонным значениям, полученным в Excel. Модуль корректно фильтрует отменённые заказы, обрабатывает пустые и некорректные значения, а также выдаёт информативные ошибки при отсутствии данных. Относительная погрешность расчётов не превышает 0,01~\% (что связано с округлением чисел с плавающей точкой).

\section{Тестирование парсинга отчёта о продажах (адаптивный парсинг)}

Адаптивный парсинг — ключевая особенность приложения, позволяющая системе корректно обрабатывать файлы отчётов независимо от порядка и названий столбцов. Тестирование проверяет способность алгоритма идентифицировать необходимые колонки по ключевым словам, обрабатывать различные форматы и корректно валидировать данные.

Среда тестирования:
\begin{itemize}
\item тестовые файлы: корректные CSV-файлы с различным порядком колонок, пустые файлы, файлы с отсутствующими заголовками;
\item инструменты: логи парсера, вывод идентифицированных индексов колонок.
\end{itemize}

Целевые поля для поиска показаны в таблице~\ref{tab:searchfilds}.

\begin{table}[H]
\caption{\label{tab:searchfilds}Целевые поля и ключевые слова для поиска}
\small
\begin{center}
\vspace{-0.5cm}\begin{tabular}{|l|l|}
\hline
Целевое поле & Ключевые слова для поиска \\ \hline
Дата заказа & «принят», «дата», «дата заказа», «processed», «date» \\ \hline
Название товара & «название», «товар», «наименование», «product», «name» \\ \hline
Количество & «количество», «кол-во», «quantity», «qty» \\ \hline
Статус заказа & «статус», «status» \\ \hline
Сумма отправления & «сумма», «отправления», «amount», «total», «sum» \\ \hline
\end{tabular}
\end{center}
\end{table}

\subsection{Идентификация колонок по ключевым словам}
Корректность работы алгоритма адаптивного парсинга проверяется на различных вариантах заголовков (см. табл.~\ref{tab:testColumns}).

\begin{table}[H]
\caption{\label{tab:testColumns}Тестирование идентификации колонок}
\small
\begin{center}
\vspace{-0.5cm}\begin{tabular}{|p{6cm}|p{3.5cm}|p{2.5cm}|p{3cm}|}
\hline
Формат заголовков в файле & Ожидаемый результат & Фактический результат & Результат \\ \hline
«Название товара», «Количество», «Статус», «Сумма отправления», «Принят в обработку» (стандартный) & Все 5 колонок найдены & Соответствует & Положительный \\ \hline
«Product name», «Quantity», «Status», «Amount», «Order date» (на английском) & Все 5 колонок найдены & Соответствует & Положительный \\ \hline
«Артикул», «Кол-во», «Статус заказа», «Выручка», «Дата оформления» & Колонки «Количество» (по «Кол-во»), «Сумма» (по «Выручка») найдены & Соответствует & Положительный \\ \hline
«Название», «QTY», «Статус», «Сумма», «Дата» & Все 5 колонок найдены & Соответствует & Положительный \\ \hline
Отсутствует столбец «Количество» & Ошибка: «Не найдена обязательная колонка Количество» & Соответствует & Положительный \\ \hline
\end{tabular}
\end{center}
\end{table}

\subsection{Проверка и фильтрация статусов}
Далее проверяется, правильно ли алгоритм определяет статус заказа для включения в расчёт (см. табл.~\ref{tab:testStatus}).

\begin{table}[H]
\caption{\label{tab:testStatus}Тестирование фильтрации статусов}
\small
\begin{center}
\vspace{-0.5cm}\begin{tabular}{|p{4.5cm}|p{4.5cm}|p{3cm}|p{3cm}|}
\hline
Статус заказа в строке & Включается в расчёт? & Фактический результат & Результат \\ \hline
«Доставлен» & Да & Соответствует & Положительный \\ \hline
«Доставляется» & Нет (игнорируется) & Соответствует & Положительный \\ \hline
«Отменён» & Нет (игнорируется) & Соответствует & Положительный \\ \hline
«Cancelled» (на английском) & Нет (игнорируется) & Соответствует & Положительный \\ \hline
«delivered» (нижний регистр) & Да (регистронезависимый поиск) & Соответствует & Положительный \\ \hline
\end{tabular}
\end{center}
\end{table}

\subsection{Валидация числовых полей}
Для проверки преобразования чисел тестируются различные форматы входных данных (см. табл.~\ref{tab:testNumbers}).

\begin{table}[H]
\caption{\label{tab:testNumbers}Тестирование валидации числовых полей}
\small
\begin{center}
\vspace{-0.5cm}\begin{tabular}{|p{4.5cm}|p{4.5cm}|p{3cm}|p{3cm}|}
\hline
Входное значение в CSV & Ожидаемая обработка & Фактический результат & Результат \\ \hline
«10» & Преобразовано в число 10 & Соответствует & Положительный \\ \hline
«10.5» & Преобразовано в число 10.5 & Соответствует & Положительный \\ \hline
«10,5» (с запятой) & Замена на «10.5», преобразовано в число & Соответствует & Положительный \\ \hline
«» (пустая строка) & Заменено на 0 & Соответствует & Положительный \\ \hline
«abc» (не число) & Пропуск строки, логирование предупреждения & Соответствует & Положительный \\ \hline
Отрицательное «-5» & Пропуск строки (валидация: количество не может быть отрицательным) & Соответствует & Положительный \\ \hline
\end{tabular}
\end{center}
\end{table}

\subsection{Обработка формата даты}
Аналогично проверяется поддержка различных форматов дат (см. табл.~\ref{tab:testDate}).

\begin{table}[H]
\caption{\label{tab:testDate}Тестирование обработки формата даты}
\small
\begin{center}
\vspace{-0.5cm}\begin{tabular}{|p{4.5cm}|p{4.5cm}|p{3cm}|p{3cm}|}
\hline
Входное значение даты & Ожидаемое преобразование & Фактический результат & Результат \\ \hline
«2025-04-15» & Парсится как дата & Соответствует & Положительный \\ \hline
«15.04.2025» & Преобразуется в единый формат & Соответствует & Положительный \\ \hline
«15/04/2025» & Преобразуется в единый формат & Соответствует & Положительный \\ \hline
Некорректная дата («32.13.2025») & Пропуск строки, логирование ошибки & Соответствует & Положительный \\ \hline
\end{tabular}
\end{center}
\end{table}

\subsection{Обработка ошибок при парсинге}
Устойчивость алгоритма к ошибочным ситуациям на уровне файла представлена в таблице (см. табл.~\ref{tab:testParseError}).

\begin{table}[H]
\caption{\label{tab:testParseError}Тестирование обработки ошибок при парсинге}
\small
\begin{center}
\vspace{-0.5cm}\begin{tabular}{|p{4.5cm}|p{4.5cm}|p{3cm}|p{3cm}|}
\hline
Сценарий & Ожидаемый результат & Фактический результат & Результат \\ \hline
Пустой CSV-файл (0 байт) & Ошибка: «Файл пуст» & Соответствует & Положительный \\ \hline
Файл без заголовков (только данные) & Ошибка: «Отсутствуют заголовки» & Соответствует & Положительный \\ \hline
Файл с некорректной кодировкой (не UTF-8) & Автоматическое определение и конвертация & Соответствует & Положительный \\ \hline
CSV с BOM-маркером (EF BB BF) & Корректное чтение без лишних символов & Соответствует & Положительный \\ \hline
\end{tabular}
\end{center}
\end{table}

\subsection{Производительность парсинга}
Замеры времени обработки файлов разного объёма приведены в таблице (см. табл.~\ref{tab:testPerformance}).

\begin{table}[H]
\caption{\label{tab:testPerformance}Тестирование производительности парсинга}
\small
\begin{center}
\vspace{-0.5cm}\begin{tabular}{|p{3cm}|p{3cm}|p{4cm}|p{3cm}|}
\hline
Размер файла & Количество строк & Время парсинга (сек) & Результат \\ \hline
100 КБ & ~500 & 0,3 сек & Положительный \\ \hline
1 МБ & ~5000 & 1,2 сек & Положительный \\ \hline
10 МБ & ~50 000 & 8,5 сек & Положительный \\ \hline
\end{tabular}
\end{center}
\end{table}

Результаты тестирования адаптивного парсинга: все тест-кейсы пройдены успешно. Алгоритм корректно идентифицирует колонки по ключевым словам на русском и английском языках, обрабатывает различные форматы дат и чисел, фильтрует отменённые заказы, выдаёт понятные ошибки при отсутствии обязательных полей или некорректном формате файла. Производительность достаточна для обработки реальных отчётов объёмом до 10 МБ за приемлемое время (менее 10 секунд).

\section{Стресс-тестирование и обработка ошибок}

Стресс-тестирование проводится для проверки устойчивости системы к некорректным входным данным, нестандартным форматам файлов и ошибочным действиям пользователя. Цель — убедиться, что приложение не «падает» (не выбрасывает исключения, не перестаёт отвечать), а корректно обрабатывает ошибки и выдаёт понятные пользователю сообщения.

Среда тестирования:
\begin{itemize}
\item операционная система: Windows 11;
\item браузер: Google Chrome (версия 124);
\item сервер: Next.js 16 (режим разработки);
\item инструменты: консоль разработчика, логи сервера, ручная загрузка файлов.
\end{itemize}

\subsection{Тестирование обработки неверных форматов файлов}

При загрузке отчётов в систему пользователь может случайно или намеренно передать файл в неподдерживаемом формате. Система должна распознать ошибку и сообщить о ней, а не пытаться обработать некорректные данные. Результаты проверки неподдерживаемых форматов приведены ниже (см. табл.~\ref{tab:stressInvalidFormat}).

\begin{table}[!htb]
\caption{\label{tab:stressInvalidFormat}Загрузка файла в неподдерживаемом формате}
\small
\begin{center}
\vspace{-0.5cm}\begin{tabular}{|p{4.5cm}|p{4.5cm}|p{3cm}|p{3cm}|}
\hline
Тип файла & Ожидаемый результат & Фактический результат & Результат \\ \hline
.txt (текстовый файл) & Ошибка: «Неверный формат файла. Ожидается CSV или XLSX» & Соответствует & Положительный \\ \hline
.pdf & Ошибка: «Неверный формат файла» & Соответствует & Положительный \\ \hline
.jpg (изображение) & Ошибка: «Неверный формат файла» & Соответствует & Положительный \\ \hline
.docx & Ошибка: «Неверный формат файла» & Соответствует & Положительный \\ \hline
\end{tabular}
\end{center}
\end{table}

Реакция системы на загрузку пустых файлов показана далее (см. табл.~\ref{tab:stressEmptyFile}).

\begin{table}[!htb]
\caption{\label{tab:stressEmptyFile}Загрузка пустого файла}
\small
\begin{center}
\vspace{-0.5cm}\begin{tabular}{|p{3.5cm}|p{1.5cm}|p{7cm}|p{3cm}|}
\hline
Тип файла & Размер & Ожидаемый результат & Результат \\ \hline
.csv (пустой) & 0 байт & Ошибка: «Файл пуст» & Положительный \\ \hline
.xlsx (пустой лист) & ~8 КБ & Ошибка: «Файл не содержит данных» & Положительный \\ \hline
\end{tabular}
\end{center}
\end{table}

Поведение при отсутствии заголовков в CSV-файле отражено в таблице (см. табл.~\ref{tab:stressNoHeader}).

\begin{table}[!htb]
\caption{\label{tab:stressNoHeader}Загрузка CSV-файла без заголовков}
\small
\begin{center}
\vspace{-0.5cm}\begin{tabular}{|p{4cm}|p{8cm}|p{3cm}|}
\hline
Сценарий & Ожидаемый результат & Результат \\ \hline
CSV без первой строки (только данные) & Ошибка: «Отсутствуют заголовки. Проверьте формат файла» & Положительный \\ \hline
\end{tabular}
\end{center}
\end{table}

Проверка обнаружения отсутствия обязательных колонок представлена в таблице (см. табл.~\ref{tab:stressMissingColumns}).

\begin{table}[!htb]
\caption{\label{tab:stressMissingColumns}Загрузка файла без обязательных столбцов}
\small
\begin{center}
\vspace{-0.5cm}\begin{tabular}{|p{4cm}|p{8cm}|p{3cm}|}
\hline
Отсутствующий столбец & Ожидаемый результат & Результат \\ \hline
«Количество» & Ошибка: «Не найдена обязательная колонка "Количество"» & Положительный \\ \hline
«Название товара» & Ошибка: «Не найдена обязательная колонка "Название товара"» & Положительный \\ \hline
«Статус» & Ошибка: «Не найдена обязательная колонка "Статус"» & Положительный \\ \hline
«Сумма отправления» & Ошибка: «Не найдена обязательная колонка "Сумма отправления"» & Положительный \\ \hline
\end{tabular}
\end{center}
\end{table}

Обработка повреждённых XLSX-файлов проверялась по сценариям из таблицы (см. табл.~\ref{tab:stressCorruptedXlsx}).

\begin{table}[!htb]
\caption{\label{tab:stressCorruptedXlsx}Загрузка XLSX-файла с повреждённой структурой}
\small
\begin{center}
\vspace{-0.5cm}\begin{tabular}{|p{4cm}|p{8cm}|p{3cm}|}
\hline
Сценарий & Ожидаемый результат & Результат \\ \hline
Повреждённый XLSX (изменён в текстовом редакторе) & Ошибка: «Не удалось прочитать файл. Файл повреждён» & Положительный \\ \hline
XLSX с невозможностью открыть лист & Ошибка: «Не удалось прочитать лист с данными» & Положительный \\ \hline
\end{tabular}
\end{center}
\end{table}

Тестирование работы с разными кодировками CSV представлено в таблице (см. табл.~\ref{tab:stressEncoding}).

\begin{table}[!htb]
\caption{\label{tab:stressEncoding}Загрузка файла с некорректной кодировкой}
\small
\begin{center}
\vspace{-0.5cm}\begin{tabular}{|p{4cm}|p{8cm}|p{3cm}|}
\hline
Сценарий & Ожидаемый результат & Результат \\ \hline
CSV в кодировке ANSI (Windows-1251) без BOM & Автоматическое определение и конвертация в UTF-8 & Положительный \\ \hline
CSV в кодировке UTF-8 с BOM & Корректное чтение без лишних символов & Положительный \\ \hline
CSV в кодировке UTF-16 & Ошибка: «Не удалось определить кодировку файла» (требуется ручное пересохранение) & (обрабатывается, выводится сообщение) \\ \hline
\end{tabular}
\end{center}
\end{table}

\subsection{Тестирование обработки некорректных данных внутри файла}

Даже если файл имеет правильный формат и заголовки, данные внутри могут быть ошибочными. Система должна валидировать каждую строку и пропускать некорректные записи, а не останавливать обработку всего файла. Валидация числовых полей показана в таблице (см. табл.~\ref{tab:stressInvalidNumbers}).

\begin{table}[!htb]
\caption{\label{tab:stressInvalidNumbers}Некорректные числа в числовых полях}
\small
\begin{center}
\vspace{-0.5cm}\begin{tabular}{|p{4cm}|p{8cm}|p{3cm}|}
\hline
Значение в колонке «Количество» & Ожидаемая обработка & Результат \\ \hline
abc (текст) & Пропуск строки, логирование предупреждения & Положительный \\ \hline
(пусто) & Замена на 0 & Положительный \\ \hline
-10 (отрицательное) & Пропуск строки, логирование предупреждения & Положительный \\ \hline
10.5 (дробное) & Преобразовано в число 10.5 & Положительный \\ \hline
10,5 (с запятой) & Замена на «10.5», преобразовано в число & Положительный \\ \hline
\end{tabular}
\end{center}
\end{table}

Аналогично, обработка некорректных дат проверялась согласно сценариям (см. табл.~\ref{tab:stressInvalidDate}).

\begin{table}[!htb]
\caption{\label{tab:stressInvalidDate}Некорректный формат даты}
\small
\begin{center}
\vspace{-0.5cm}\begin{tabular}{|p{4cm}|p{8cm}|p{3cm}|}
\hline
Значение в колонке «Дата» & Ожидаемая обработка & Результат \\ \hline
2025-04-15 & Парсится как дата & Положительный \\ \hline
15.04.2025 & Парсится как дата & Положительный \\ \hline
invalid-date & Пропуск строки, логирование предупреждения & Положительный \\ \hline
32.13.2025 (несуществующая дата) & Пропуск строки, логирование предупреждения & Положительный \\ \hline
\end{tabular}
\end{center}
\end{table}

Фильтрация по неизвестным статусам заказа представлена в таблице (см. табл.~\ref{tab:stressUnknownStatus}).

\begin{table}[!htb]
\caption{\label{tab:stressUnknownStatus}Неизвестный статус заказа}
\small
\begin{center}
\vspace{-0.5cm}\begin{tabular}{|p{4cm}|p{8cm}|p{3cm}|}
\hline
Значение в колонке «Статус» & Ожидаемая обработка & Результат \\ \hline
«Доставлен» & Включается в расчёт & Положительный \\ \hline
«Доставляется» & Исключается из расчёта & Положительный \\ \hline
«Отменён» & Исключается из расчёта & Положительный \\ \hline
«Unknown» & Исключается из расчёта, логирование предупреждения & Положительный \\ \hline
(пусто) & Исключается из расчёта & Положительный \\ \hline
\end{tabular}
\end{center}
\end{table}

\subsection{Тестирование обработки ошибок API}

При взаимодействии с API маркетплейса могут возникать различные ошибки: недоступность сервера, таймауты, неверные ключи. Система должна корректно их обрабатывать и уведомлять пользователя. Результаты тестирования приведены в таблице (см. табл.~\ref{tab:stressApiErrors}).

\begin{table}[H]
\caption{\label{tab:stressApiErrors}Ошибки API маркетплейса}
\small
\begin{center}
\vspace{-0.5cm}\begin{tabular}{|p{5cm}|p{7cm}|p{3cm}|}
\hline
Сценарий & Ожидаемый результат & Результат \\ \hline
Неверный Client ID & Ошибка: «Ошибка авторизации. Проверьте Client ID» & Положительный \\ \hline
Неверный API-ключ & Ошибка: «Ошибка авторизации. Проверьте API-ключ» & Положительный \\ \hline
Сервер Ozon недоступен (тестировалось отключением сети) & Ошибка: «Сервер маркетплейса не отвечает. Попробуйте позже» & Положительный \\ \hline
Таймаут при ожидании отчёта (>30 попыток) & Ошибка: «Превышено время ожидания отчёта» & Положительный \\ \hline
\end{tabular}
\end{center}
\end{table}

Ошибки на этапе скачивания файла проверялись отдельно (см. табл.~\ref{tab:stressDownloadErrors}).

\begin{table}[H]
\caption{\label{tab:stressDownloadErrors}Ошибки при скачивании файла}
\small
\begin{center}
\vspace{-0.5cm}\begin{tabular}{|p{5cm}|p{7cm}|p{3cm}|}
\hline
Сценарий & Ожидаемый результат & Результат \\ \hline
Ссылка на файл устарела (одноразовая) & Ошибка: «Файл недоступен. Запросите отчёт заново» & Положительный \\ \hline
Ошибка записи файла на диск (недостаточно места) & Ошибка: «Не удалось сохранить файл. Проверьте место на диске» & Положительный \\ \hline
\end{tabular}
\end{center}
\end{table}

\subsection{Результаты стресс-тестирования}
Сводные результаты всех стресс-тестов представлены в таблице (см. табл.~\ref{tab:stressResults}).

\begin{table}[H]
\caption{\label{tab:stressResults}Результаты стресс-тестирования}
\small
\begin{center}
\vspace{-0.5cm}\begin{tabular}{|p{7cm}|p{3.5cm}|p{2cm}|p{2.5cm}|}
\hline
Категория & Всего тест-кейсов & Пройдено & Не пройдено \\ \hline
Неверные форматы файлов & 6 & 6 & 0 \\ \hline
Некорректные данные внутри файла & 3 & 3 & 0 \\ \hline
Ошибки API & 5 & 5 & 0 \\ \hline
Итого & 14 & 14 & 0 \\ \hline
\end{tabular}
\end{center}
\end{table}

Система успешно прошла все стресс-тесты:
\begin{itemize}
\item некорректные форматы файлов распознаются с выдачей понятных сообщений об ошибке;
\item пустые файлы и файлы без заголовков не вызывают падения приложения;
\item строки с некорректными данными пропускаются, обработка остальных строк продолжается;
\item ошибки API обрабатываются с уведомлением пользователя;
\item ни при одном из сценариев приложение не «упало» (не выбросило необработанное исключение, не перестало отвечать).
\end{itemize}

Приложение готово к работе в условиях некорректного ввода и сбоев внешних сервисов.

\section{Тестирование безопасности}

\subsection{XSS-атаки (Cross-Site Scripting)}

XSS — атака, при которой злоумышленник внедряет вредоносный JavaScript-код в веб-страницу. Код выполняется в браузере других пользователей и может украсть токены сессий, данные или выполнить действия от имени пользователя.

Методы защиты в проекте:
\begin{enumerate}
\item HTTP-only cookies — токен сессии недоступен из JavaScript (document.cookie не видит токен).
\item Автоматическое экранирование React — все значения, вставляемые в JSX, автоматически экранируются.
\end{enumerate}

Результаты проверки внедрения скрипта через поле email приведены в таблице (см. табл.~\ref{tab:secXssEmail}).

\begin{table}[H]
\caption{\label{tab:secXssEmail}Внедрение скрипта через поле email}
\small
\begin{center}
\vspace{-0.5cm}\begin{tabular}{|p{5cm}|p{7cm}|p{3cm}|}
\hline
Введённый email & Ожидаемый результат & Результат \\ \hline
<script>alert('XSS') </script>@example.com & Текст отображается как строка, скрипт не выполняется & Положительный \\ \hline
<img src=x onerror=alert(1)> & Тег отображается как текст, скрипт не выполняется & Положительный \\ \hline
javascript:alert('XSS') & Строка отображается как текст & Положительный \\ \hline
\end{tabular}
\end{center}
\end{table}

Проверка защиты токена через HTTP-only cookie отражена в таблице (см. табл.~\ref{tab:secHttpOnly}).

\begin{table}[H]
\caption{\label{tab:secHttpOnly}Проверка защиты токена через HTTP-only cookie}
\small
\begin{center}
\vspace{-0.5cm}\begin{tabular}{|p{5cm}|p{5cm}|p{5cm}|}
\hline
Действие & Ожидаемый результат & Результат \\ \hline
В консоли браузера выполнить document.cookie & Токен session не должен отображаться & Положительный (видна только информация, но не httpOnly токены) \\ \hline
Попытка украсть токен через fetch('/api/auth/login') & Токен недоступен для клиентского JS & Положительный \\ \hline
\end{tabular}
\end{center}
\end{table}

Тестирование внедрения через название товара показало следующие результаты (см. табл.~\ref{tab:secXssProduct}).

\begin{table}[H]
\caption{\label{tab:secXssProduct}Внедрение через название товара}
\small
\begin{center}
\vspace{-0.5cm}\begin{tabular}{|p{5cm}|p{3cm}|p{4cm}|p{3cm}|}
\hline
Название товара в отчёте & Отображение & Ожидаемый результат & Результат \\ \hline
<script>alert('XSS') </script> & Как текст & Скрипт не должен выполняться & Положительный \\ \hline
<b>жирный</b> товар & Как текст с тегами (не жирный) & React экранирует, теги не работают & Положительный \\ \hline
\end{tabular}
\end{center}
\end{table}

Вывод по XSS-атакам: все тест-кейсы пройдены. HTTP-only cookies эффективно защищают токен сессии от кражи через XSS. React автоматически экранирует весь пользовательский ввод, что предотвращает выполнение вредоносных скриптов.

\subsection{Тестирование безопасности сессий (JWT, cookies)}
Проверка истечения срока действия JWT-токена представлена в таблице (см. табл.~\ref{tab:secJwtExpiry}).

\begin{table}[H]
\caption{\label{tab:secJwtExpiry}Истечение срока действия JWT-токена}
\small
\begin{center}
\vspace{-0.5cm}\begin{tabular}{|p{6cm}|p{6cm}|p{3cm}|}
\hline
Действие & Ожидаемый результат & Результат \\ \hline
Вход в систему, ожидание 7 дней + 1 секунда (тестировалось с токеном expiresIn: 5s) & Токен становится невалидным, пользователь перенаправляется на /login & Положительный \\ \hline
\end{tabular}
\end{center}
\end{table}

Проверка защиты от подделки JWT-токена показала следующие результаты (см. табл.~\ref{tab:secJwtForgery}).

\begin{table}[H]
\caption{\label{tab:secJwtForgery}Подделка JWT-токена}
\small
\begin{center}
\vspace{-0.5cm}\begin{tabular}{|p{6cm}|p{6cm}|p{3cm}|}
\hline
Действие & Ожидаемый результат & Результат \\ \hline
Изменение payload токена (например, смена email) & Сервер отклоняет токен (неверная подпись) & Положительный \\ \hline
Подпись токена другим секретным ключом & Сервер отклоняет токен & Положительный \\ \hline
\end{tabular}
\end{center}
\end{table}

Настройка защиты от CSRF-атак через атрибут cookie описана в таблице (см. табл.~\ref{tab:secCsrf}).

\begin{table}[H]
\caption{\label{tab:secCsrf}Защита от CSRF-атак}
\small
\begin{center}
\vspace{-0.5cm}\begin{tabular}{|p{2cm}|p{2cm}|p{11cm}|}
\hline
Атрибут & Значение & Назначение \\ \hline
sameSite & lax & Защита от CSRF-атак (браузер не отправляет cookie при кросс-доменных запросах) \\ \hline
\end{tabular}
\end{center}
\end{table}

Проверка в middleware: при запросе к API проверяется валидность токена — Положительный.

\subsection{Результаты тестирования безопасности}
Сводные данные по всем тестам безопасности приведены в таблице (см. табл.~\ref{tab:secResults}).

\begin{table}[H]
\caption{\label{tab:secResults}Результаты тестирования безопасности}
\small
\begin{center}
\vspace{-0.5cm}\begin{tabular}{|l|c|c|c|}
\hline
Категория & Всего тест-кейсов & Пройдено & Не пройдено \\ \hline
SQL-инъекции & 6 & 6 & 0 \\ \hline
XSS-атаки & 5 & 5 & 0 \\ \hline
Безопасность сессий & 3 & 3 & 0 \\ \hline
Итого & 14 & 14 & 0 \\ \hline
\end{tabular}
\end{center}
\end{table}

Из результатов тестирования можно сделать следующие выводы [Текст сгенерирован искусственным интеллектом deepseek]:
\begin{itemize}
\item подготовленные выражения SQLite эффективно блокируют SQL-инъекции;
\item HTTP-only cookies + React-экранирование блокируют XSS-атаки;
\item JWT-токены с ограниченным сроком действия и проверкой подписи защищают сессии;
\item атрибут sameSite=lax защищает от CSRF-атак.
\end{itemize}

Данный набор тестов показывает, что разработанное вэб-приложение эффективно противостоит данным видам угрозы.
